% Conclusions.  A fragment: include from the master with % Conclusions.  A fragment: include from the master with % Conclusions.  A fragment: include from the master with % Conclusions.  A fragment: include from the master with \input{conclusions}.
% Requires nothing beyond booktabs/amsmath/siunitx, which the master loads.

\section{Conclusions}
\subsection{What was established}

The thesis set out to compute a defensible pump-and-treat design for a
contaminated aquifer under conductivity uncertainty, and to be explicit about
what makes such a design defensible. Three results carry the argument.

\paragraph{The surrogate question.} On the one-dimensional advection--dispersion
sandbox, with the Ogata--Banks solution as the reference, a physics-informed
network reaches a median full-grid relative error of $1.61\times10^{-2}$ over
five seeds, with individual seeds spanning $1.17\times10^{-2}$ to
$8.14\times10^{-2}$; a coordinate-augmented Fourier neural operator reaches a
median of $1.99\times10^{-2}$ at 250 epochs over five seeds, $11\times$ better
than the mean-field predictor, where its best seed alone would have claimed
$1.46\times10^{-2}$. The same operator, in an earlier configuration, scored
$0.973$ --- worse than outputting zero, because relative error is normalized by
the target and the architecture could only produce a constant field. Both facts
are only visible against a baseline and a seed sweep.

\paragraph{The optimization question.} On the $200\times120$\,m, $160\times96$
benchmark with $1\,529$ Karhunen--Lo\`eve modes and a $5\%$ chance constraint,
two independent runs that differ only in floating-point precision converge to
policies $1.2\%$ apart in penalized cost on different well subsets
($J = 428\,592$ and $433\,945$), each sitting on the constraint
($\mathrm{CVaR}_{5\%}\approx 0$) as the solution of a chance-constrained program
must, in $1.67$\,h and $4.14$\,h of one rented GPU. Batching an entire compass
sweep is worth $47\times$ over the NumPy reference; compiling the Krylov
iteration into a single call is worth a further $2.2$--$2.8\times$ at the larger
geometry, with the objective unchanged to $10^{-7}$.

\paragraph{The certification question.} A policy is reported only after being
re-solved on scenarios no run has seen, on the independent NumPy reference. The
accepted policy gives $391/400$ compliant scenarios on a fresh draw,
reliability $0.978$ with a Wilson interval of $[0.958, 0.988]$, clearing the
$0.94$ gate; its cost reproduces to $0.02\%$. The same protocol rejects a policy
optimized on sixteen scenarios: it reported reliability $1.000$ on its own
sample and achieves $0.912$ $[0.881, 0.936]$ on four hundred unseen ones.

\subsection{The methodological result}

The finding that generalizes beyond this problem is a statement about which
statistics of a stochastic program transport from the optimization sample to the
aquifer.

For the accepted policy on three disjoint draws of $400$ scenarios, the cost
reproduces to $0.02\%$ and the reliability to $0.003$
($0.978$, $0.975$, $0.978$), while the CVaR tail takes the values $+0.616$,
$+0.254$ and $+0.334$. Those three are not distinguishable: their spread is
smaller than the standard error of any one of them ($\pm0.21$ to $\pm0.47$),
because a $5\%$ tail of $400$ draws averages twenty of the most variable
realizations in the set. The reported in-sample value, $+0.005$, lies outside all
three intervals --- it is a fourth sample, not a better estimate.

The practical consequences, which are cheap to adopt and are the recommendation
of this thesis, are fourfold. Certify a policy with \emph{reliability and its
interval}, not with a tail value. Quote a tail \emph{with its standard error}, and
only as evidence about where the policy sits relative to the constraint. Do not
rank policies by tail value at sample sizes that are affordable: at $M=400$ the
tail cannot separate policies that reliability separates easily. And re-evaluate
on draws that no part of the optimization has touched, because scenarios that
continue the training seed reproduce the run rather than test it.

\subsection{Limitations of the present evidence}

\begin{itemize}
  \item \textbf{One two-dimensional geometry.} All 2D conclusions concern a single
        $200\times120$\,m configuration, nine wells and one source. The
        observation that well rates are seed-insensitive --- two independently
        seeded searches followed a digit-for-digit identical trajectory --- is a
        property of this cost model, not a general law, and it means the reported
        policy is one of several near-equivalent designs rather than a unique
        optimum.
  \item \textbf{Mixed precision at the margin.} The accepted policy was found in
        mixed float32 (float32 operator, float64 Krylov recurrence). That mode is
        not unconditionally safe: two of three fully float32 runs aborted on a
        non-finite transport residual, which is why the search now tolerates a
        single rejected candidate but not an entire diverging sweep. The
        float32/float64 policy gap of $1.2\%$ is within the noise that fine
        compass steps introduce.
  \item \textbf{Tail resolution.} At $M=400$ the tail statistic is not resolvable
        to better than a few tenths. Establishing the tail of the distribution to
        the precision at which reliability is known would require thousands of
        scenarios, i.e.\ an order of magnitude more transport solves than the
        optimization itself.
  \item \textbf{The three-dimensional model is verified but not certified, and
        is coarse.} Seven gates are asserted on the solver --- exact seepage
        velocity, transversally invariant head, a Thiem shape to $8.1\times
        10^{-3}$, agreement with the 2D code under extrusion to $10^{-13}$, the
        advection front speed to $1.2\%$, a breakthrough variance consistent
        with $\alpha_L v$, and a mass ledger closing to $10^{-15}$ --- but there is
        no analytical \emph{shape} gate for 3D transport, deliberately: the
        Ogata--Banks solution is a continuous Dirichlet step, and against a finite
        inlet it differs by a truncation error that shrinks with time ($0.32\to
        0.13$), so it would test the upwinding both codes share rather than the
        dimensional extension. What the 3D results lack is a \emph{certificate}:
        the paper preset's interval is $[0.922,0.988]$ and the $8\,000$-cell
        policy's is $[0.872,0.963]$, neither clearing the $0.94$ gate, and the
        4\,096-cell geometry is an order of magnitude coarser than the 2D
        campaign with a single log-normal conductivity family rather than a
        layered geology.
  \item \textbf{Timings are single-machine.} Reported GPU costs come from rented
        T4 and A100 instances with the batch sizes stated; the fixed-cost model
        behind the larger projections is a two-point fit that was not
        reproducible to better than about $15\%$.
\end{itemize}

\subsection{Work that would remove these limitations}

\begin{enumerate}
  \item \textbf{A certified 3D policy, and an affordable way to compute one.}
        Certifying at the 2D standard means $M_{\text{test}}\approx400$ and
        $M_{\text{train}}\approx100$ --- about $18$\,h per two-start search on
        one CPU at $8\,000$ cells --- so the work is as much computational as
        statistical: the batched torch backend on the $32\,768$-cell preset, and
        a search that uses the adjoint instead of a thousand forward solves.
  \item \textbf{A 2.5D axisymmetric bridge} between the 2D and 3D models, which
        would give a second independent route to the same answer and make the
        dimensional reduction a check instead of an assumption.
  \item \textbf{Operator learning in 3D} as an input--output map between
        conductivity realizations and compliance-point histories, with the same
        seed-sweep and resolution-transfer protocol applied in 1D, so that the
        surrogate claim is either established or refuted at the resolution the
        optimization actually uses.
  \item \textbf{Tail estimation at scale.} If a tail statement is wanted, it has
        to be paid for: $M$ of a few thousand, drawn from a frozen policy, with
        the standard error reported as a matter of course. Until then the honest
        claim is about reliability.
  \item \textbf{Adaptive design, solved by progressive hedging.} The formulation
        is written down and checked --- first-stage rates, recourse rates that may
        depend on what has been observed at $t_{\text{obs}}$, nonanticipativity on
        the first stage only, and a value of the stochastic solution and of perfect
        information that a here-and-now model cannot report --- and the
        decomposition on the one-stage problem is explicitly degenerate (a
        nine-dimensional consensus block has nothing to decompose). Running it at
        benchmark scale is the experiment that would say whether observing an
        aquifer for a year is cheaper than installing more wells, which is the
        question the 3D cost wall makes urgent.
  \item \textbf{Removing the per-step kernel launches.} The measured $2.2$--$2.8\times$
        came from compiling the iteration body; the sequential dependence of
        implicit time stepping remains, and a formulation that solves the whole
        time history at once would remove the remaining synchronizations. Any
        such change must still pass the mass-balance and field-parity gates.
\end{enumerate}

The wider lesson is that the two questions the thesis opened with --- is the
forward model right, and is the optimizer's evidence about the aquifer or about
its own sample --- are answered by the same kind of device: an independent
computation that can disagree. The cost of each device is small relative to the
optimization it audits, and in this work every one of them changed a conclusion.
.
% Requires nothing beyond booktabs/amsmath/siunitx, which the master loads.

\section{Conclusions}
\subsection{What was established}

The thesis set out to compute a defensible pump-and-treat design for a
contaminated aquifer under conductivity uncertainty, and to be explicit about
what makes such a design defensible. Three results carry the argument.

\paragraph{The surrogate question.} On the one-dimensional advection--dispersion
sandbox, with the Ogata--Banks solution as the reference, a physics-informed
network reaches a median full-grid relative error of $1.61\times10^{-2}$ over
five seeds, with individual seeds spanning $1.17\times10^{-2}$ to
$8.14\times10^{-2}$; a coordinate-augmented Fourier neural operator reaches a
median of $1.99\times10^{-2}$ at 250 epochs over five seeds, $11\times$ better
than the mean-field predictor, where its best seed alone would have claimed
$1.46\times10^{-2}$. The same operator, in an earlier configuration, scored
$0.973$ --- worse than outputting zero, because relative error is normalized by
the target and the architecture could only produce a constant field. Both facts
are only visible against a baseline and a seed sweep.

\paragraph{The optimization question.} On the $200\times120$\,m, $160\times96$
benchmark with $1\,529$ Karhunen--Lo\`eve modes and a $5\%$ chance constraint,
two independent runs that differ only in floating-point precision converge to
policies $1.2\%$ apart in penalized cost on different well subsets
($J = 428\,592$ and $433\,945$), each sitting on the constraint
($\mathrm{CVaR}_{5\%}\approx 0$) as the solution of a chance-constrained program
must, in $1.67$\,h and $4.14$\,h of one rented GPU. Batching an entire compass
sweep is worth $47\times$ over the NumPy reference; compiling the Krylov
iteration into a single call is worth a further $2.2$--$2.8\times$ at the larger
geometry, with the objective unchanged to $10^{-7}$.

\paragraph{The certification question.} A policy is reported only after being
re-solved on scenarios no run has seen, on the independent NumPy reference. The
accepted policy gives $391/400$ compliant scenarios on a fresh draw,
reliability $0.978$ with a Wilson interval of $[0.958, 0.988]$, clearing the
$0.94$ gate; its cost reproduces to $0.02\%$. The same protocol rejects a policy
optimized on sixteen scenarios: it reported reliability $1.000$ on its own
sample and achieves $0.912$ $[0.881, 0.936]$ on four hundred unseen ones.

\subsection{The methodological result}

The finding that generalizes beyond this problem is a statement about which
statistics of a stochastic program transport from the optimization sample to the
aquifer.

For the accepted policy on three disjoint draws of $400$ scenarios, the cost
reproduces to $0.02\%$ and the reliability to $0.003$
($0.978$, $0.975$, $0.978$), while the CVaR tail takes the values $+0.616$,
$+0.254$ and $+0.334$. Those three are not distinguishable: their spread is
smaller than the standard error of any one of them ($\pm0.21$ to $\pm0.47$),
because a $5\%$ tail of $400$ draws averages twenty of the most variable
realizations in the set. The reported in-sample value, $+0.005$, lies outside all
three intervals --- it is a fourth sample, not a better estimate.

The practical consequences, which are cheap to adopt and are the recommendation
of this thesis, are fourfold. Certify a policy with \emph{reliability and its
interval}, not with a tail value. Quote a tail \emph{with its standard error}, and
only as evidence about where the policy sits relative to the constraint. Do not
rank policies by tail value at sample sizes that are affordable: at $M=400$ the
tail cannot separate policies that reliability separates easily. And re-evaluate
on draws that no part of the optimization has touched, because scenarios that
continue the training seed reproduce the run rather than test it.

\subsection{Limitations of the present evidence}

\begin{itemize}
  \item \textbf{One two-dimensional geometry.} All 2D conclusions concern a single
        $200\times120$\,m configuration, nine wells and one source. The
        observation that well rates are seed-insensitive --- two independently
        seeded searches followed a digit-for-digit identical trajectory --- is a
        property of this cost model, not a general law, and it means the reported
        policy is one of several near-equivalent designs rather than a unique
        optimum.
  \item \textbf{Mixed precision at the margin.} The accepted policy was found in
        mixed float32 (float32 operator, float64 Krylov recurrence). That mode is
        not unconditionally safe: two of three fully float32 runs aborted on a
        non-finite transport residual, which is why the search now tolerates a
        single rejected candidate but not an entire diverging sweep. The
        float32/float64 policy gap of $1.2\%$ is within the noise that fine
        compass steps introduce.
  \item \textbf{Tail resolution.} At $M=400$ the tail statistic is not resolvable
        to better than a few tenths. Establishing the tail of the distribution to
        the precision at which reliability is known would require thousands of
        scenarios, i.e.\ an order of magnitude more transport solves than the
        optimization itself.
  \item \textbf{The three-dimensional model is verified but not certified, and
        is coarse.} Seven gates are asserted on the solver --- exact seepage
        velocity, transversally invariant head, a Thiem shape to $8.1\times
        10^{-3}$, agreement with the 2D code under extrusion to $10^{-13}$, the
        advection front speed to $1.2\%$, a breakthrough variance consistent
        with $\alpha_L v$, and a mass ledger closing to $10^{-15}$ --- but there is
        no analytical \emph{shape} gate for 3D transport, deliberately: the
        Ogata--Banks solution is a continuous Dirichlet step, and against a finite
        inlet it differs by a truncation error that shrinks with time ($0.32\to
        0.13$), so it would test the upwinding both codes share rather than the
        dimensional extension. What the 3D results lack is a \emph{certificate}:
        the paper preset's interval is $[0.922,0.988]$ and the $8\,000$-cell
        policy's is $[0.872,0.963]$, neither clearing the $0.94$ gate, and the
        4\,096-cell geometry is an order of magnitude coarser than the 2D
        campaign with a single log-normal conductivity family rather than a
        layered geology.
  \item \textbf{Timings are single-machine.} Reported GPU costs come from rented
        T4 and A100 instances with the batch sizes stated; the fixed-cost model
        behind the larger projections is a two-point fit that was not
        reproducible to better than about $15\%$.
\end{itemize}

\subsection{Work that would remove these limitations}

\begin{enumerate}
  \item \textbf{A certified 3D policy, and an affordable way to compute one.}
        Certifying at the 2D standard means $M_{\text{test}}\approx400$ and
        $M_{\text{train}}\approx100$ --- about $18$\,h per two-start search on
        one CPU at $8\,000$ cells --- so the work is as much computational as
        statistical: the batched torch backend on the $32\,768$-cell preset, and
        a search that uses the adjoint instead of a thousand forward solves.
  \item \textbf{A 2.5D axisymmetric bridge} between the 2D and 3D models, which
        would give a second independent route to the same answer and make the
        dimensional reduction a check instead of an assumption.
  \item \textbf{Operator learning in 3D} as an input--output map between
        conductivity realizations and compliance-point histories, with the same
        seed-sweep and resolution-transfer protocol applied in 1D, so that the
        surrogate claim is either established or refuted at the resolution the
        optimization actually uses.
  \item \textbf{Tail estimation at scale.} If a tail statement is wanted, it has
        to be paid for: $M$ of a few thousand, drawn from a frozen policy, with
        the standard error reported as a matter of course. Until then the honest
        claim is about reliability.
  \item \textbf{Adaptive design, solved by progressive hedging.} The formulation
        is written down and checked --- first-stage rates, recourse rates that may
        depend on what has been observed at $t_{\text{obs}}$, nonanticipativity on
        the first stage only, and a value of the stochastic solution and of perfect
        information that a here-and-now model cannot report --- and the
        decomposition on the one-stage problem is explicitly degenerate (a
        nine-dimensional consensus block has nothing to decompose). Running it at
        benchmark scale is the experiment that would say whether observing an
        aquifer for a year is cheaper than installing more wells, which is the
        question the 3D cost wall makes urgent.
  \item \textbf{Removing the per-step kernel launches.} The measured $2.2$--$2.8\times$
        came from compiling the iteration body; the sequential dependence of
        implicit time stepping remains, and a formulation that solves the whole
        time history at once would remove the remaining synchronizations. Any
        such change must still pass the mass-balance and field-parity gates.
\end{enumerate}

The wider lesson is that the two questions the thesis opened with --- is the
forward model right, and is the optimizer's evidence about the aquifer or about
its own sample --- are answered by the same kind of device: an independent
computation that can disagree. The cost of each device is small relative to the
optimization it audits, and in this work every one of them changed a conclusion.
.
% Requires nothing beyond booktabs/amsmath/siunitx, which the master loads.

\section{Conclusions}
\subsection{What was established}

The thesis set out to compute a defensible pump-and-treat design for a
contaminated aquifer under conductivity uncertainty, and to be explicit about
what makes such a design defensible. Three results carry the argument.

\paragraph{The surrogate question.} On the one-dimensional advection--dispersion
sandbox, with the Ogata--Banks solution as the reference, a physics-informed
network reaches a median full-grid relative error of $1.61\times10^{-2}$ over
five seeds, with individual seeds spanning $1.17\times10^{-2}$ to
$8.14\times10^{-2}$; a coordinate-augmented Fourier neural operator reaches a
median of $1.99\times10^{-2}$ at 250 epochs over five seeds, $11\times$ better
than the mean-field predictor, where its best seed alone would have claimed
$1.46\times10^{-2}$. The same operator, in an earlier configuration, scored
$0.973$ --- worse than outputting zero, because relative error is normalized by
the target and the architecture could only produce a constant field. Both facts
are only visible against a baseline and a seed sweep.

\paragraph{The optimization question.} On the $200\times120$\,m, $160\times96$
benchmark with $1\,529$ Karhunen--Lo\`eve modes and a $5\%$ chance constraint,
two independent runs that differ only in floating-point precision converge to
policies $1.2\%$ apart in penalized cost on different well subsets
($J = 428\,592$ and $433\,945$), each sitting on the constraint
($\mathrm{CVaR}_{5\%}\approx 0$) as the solution of a chance-constrained program
must, in $1.67$\,h and $4.14$\,h of one rented GPU. Batching an entire compass
sweep is worth $47\times$ over the NumPy reference; compiling the Krylov
iteration into a single call is worth a further $2.2$--$2.8\times$ at the larger
geometry, with the objective unchanged to $10^{-7}$.

\paragraph{The certification question.} A policy is reported only after being
re-solved on scenarios no run has seen, on the independent NumPy reference. The
accepted policy gives $391/400$ compliant scenarios on a fresh draw,
reliability $0.978$ with a Wilson interval of $[0.958, 0.988]$, clearing the
$0.94$ gate; its cost reproduces to $0.02\%$. The same protocol rejects a policy
optimized on sixteen scenarios: it reported reliability $1.000$ on its own
sample and achieves $0.912$ $[0.881, 0.936]$ on four hundred unseen ones.

\subsection{The methodological result}

The finding that generalizes beyond this problem is a statement about which
statistics of a stochastic program transport from the optimization sample to the
aquifer.

For the accepted policy on three disjoint draws of $400$ scenarios, the cost
reproduces to $0.02\%$ and the reliability to $0.003$
($0.978$, $0.975$, $0.978$), while the CVaR tail takes the values $+0.616$,
$+0.254$ and $+0.334$. Those three are not distinguishable: their spread is
smaller than the standard error of any one of them ($\pm0.21$ to $\pm0.47$),
because a $5\%$ tail of $400$ draws averages twenty of the most variable
realizations in the set. The reported in-sample value, $+0.005$, lies outside all
three intervals --- it is a fourth sample, not a better estimate.

The practical consequences, which are cheap to adopt and are the recommendation
of this thesis, are fourfold. Certify a policy with \emph{reliability and its
interval}, not with a tail value. Quote a tail \emph{with its standard error}, and
only as evidence about where the policy sits relative to the constraint. Do not
rank policies by tail value at sample sizes that are affordable: at $M=400$ the
tail cannot separate policies that reliability separates easily. And re-evaluate
on draws that no part of the optimization has touched, because scenarios that
continue the training seed reproduce the run rather than test it.

\subsection{Limitations of the present evidence}

\begin{itemize}
  \item \textbf{One two-dimensional geometry.} All 2D conclusions concern a single
        $200\times120$\,m configuration, nine wells and one source. The
        observation that well rates are seed-insensitive --- two independently
        seeded searches followed a digit-for-digit identical trajectory --- is a
        property of this cost model, not a general law, and it means the reported
        policy is one of several near-equivalent designs rather than a unique
        optimum.
  \item \textbf{Mixed precision at the margin.} The accepted policy was found in
        mixed float32 (float32 operator, float64 Krylov recurrence). That mode is
        not unconditionally safe: two of three fully float32 runs aborted on a
        non-finite transport residual, which is why the search now tolerates a
        single rejected candidate but not an entire diverging sweep. The
        float32/float64 policy gap of $1.2\%$ is within the noise that fine
        compass steps introduce.
  \item \textbf{Tail resolution.} At $M=400$ the tail statistic is not resolvable
        to better than a few tenths. Establishing the tail of the distribution to
        the precision at which reliability is known would require thousands of
        scenarios, i.e.\ an order of magnitude more transport solves than the
        optimization itself.
  \item \textbf{The three-dimensional model is verified but not certified, and
        is coarse.} Seven gates are asserted on the solver --- exact seepage
        velocity, transversally invariant head, a Thiem shape to $8.1\times
        10^{-3}$, agreement with the 2D code under extrusion to $10^{-13}$, the
        advection front speed to $1.2\%$, a breakthrough variance consistent
        with $\alpha_L v$, and a mass ledger closing to $10^{-15}$ --- but there is
        no analytical \emph{shape} gate for 3D transport, deliberately: the
        Ogata--Banks solution is a continuous Dirichlet step, and against a finite
        inlet it differs by a truncation error that shrinks with time ($0.32\to
        0.13$), so it would test the upwinding both codes share rather than the
        dimensional extension. What the 3D results lack is a \emph{certificate}:
        the paper preset's interval is $[0.922,0.988]$ and the $8\,000$-cell
        policy's is $[0.872,0.963]$, neither clearing the $0.94$ gate, and the
        4\,096-cell geometry is an order of magnitude coarser than the 2D
        campaign with a single log-normal conductivity family rather than a
        layered geology.
  \item \textbf{Timings are single-machine.} Reported GPU costs come from rented
        T4 and A100 instances with the batch sizes stated; the fixed-cost model
        behind the larger projections is a two-point fit that was not
        reproducible to better than about $15\%$.
\end{itemize}

\subsection{Work that would remove these limitations}

\begin{enumerate}
  \item \textbf{A certified 3D policy, and an affordable way to compute one.}
        Certifying at the 2D standard means $M_{\text{test}}\approx400$ and
        $M_{\text{train}}\approx100$ --- about $18$\,h per two-start search on
        one CPU at $8\,000$ cells --- so the work is as much computational as
        statistical: the batched torch backend on the $32\,768$-cell preset, and
        a search that uses the adjoint instead of a thousand forward solves.
  \item \textbf{A 2.5D axisymmetric bridge} between the 2D and 3D models, which
        would give a second independent route to the same answer and make the
        dimensional reduction a check instead of an assumption.
  \item \textbf{Operator learning in 3D} as an input--output map between
        conductivity realizations and compliance-point histories, with the same
        seed-sweep and resolution-transfer protocol applied in 1D, so that the
        surrogate claim is either established or refuted at the resolution the
        optimization actually uses.
  \item \textbf{Tail estimation at scale.} If a tail statement is wanted, it has
        to be paid for: $M$ of a few thousand, drawn from a frozen policy, with
        the standard error reported as a matter of course. Until then the honest
        claim is about reliability.
  \item \textbf{Adaptive design, solved by progressive hedging.} The formulation
        is written down and checked --- first-stage rates, recourse rates that may
        depend on what has been observed at $t_{\text{obs}}$, nonanticipativity on
        the first stage only, and a value of the stochastic solution and of perfect
        information that a here-and-now model cannot report --- and the
        decomposition on the one-stage problem is explicitly degenerate (a
        nine-dimensional consensus block has nothing to decompose). Running it at
        benchmark scale is the experiment that would say whether observing an
        aquifer for a year is cheaper than installing more wells, which is the
        question the 3D cost wall makes urgent.
  \item \textbf{Removing the per-step kernel launches.} The measured $2.2$--$2.8\times$
        came from compiling the iteration body; the sequential dependence of
        implicit time stepping remains, and a formulation that solves the whole
        time history at once would remove the remaining synchronizations. Any
        such change must still pass the mass-balance and field-parity gates.
\end{enumerate}

The wider lesson is that the two questions the thesis opened with --- is the
forward model right, and is the optimizer's evidence about the aquifer or about
its own sample --- are answered by the same kind of device: an independent
computation that can disagree. The cost of each device is small relative to the
optimization it audits, and in this work every one of them changed a conclusion.
.
% Requires nothing beyond booktabs/amsmath/siunitx, which the master loads.

\section{Conclusions}
\subsection{What was established}

The thesis set out to compute a defensible pump-and-treat design for a
contaminated aquifer under conductivity uncertainty, and to be explicit about
what makes such a design defensible. Three results carry the argument.

\paragraph{The surrogate question.} On the one-dimensional advection--dispersion
sandbox, with the Ogata--Banks solution as the reference, a physics-informed
network reaches a median full-grid relative error of $1.61\times10^{-2}$ over
five seeds, with individual seeds spanning $1.17\times10^{-2}$ to
$8.14\times10^{-2}$; a coordinate-augmented Fourier neural operator reaches a
median of $1.99\times10^{-2}$ at 250 epochs over five seeds, $11\times$ better
than the mean-field predictor, where its best seed alone would have claimed
$1.46\times10^{-2}$. The same operator, in an earlier configuration, scored
$0.973$ --- worse than outputting zero, because relative error is normalized by
the target and the architecture could only produce a constant field. Both facts
are only visible against a baseline and a seed sweep.

\paragraph{The optimization question.} On the $200\times120$\,m, $160\times96$
benchmark with $1\,529$ Karhunen--Lo\`eve modes and a $5\%$ chance constraint,
two independent runs that differ only in floating-point precision converge to
policies $1.2\%$ apart in penalized cost on different well subsets
($J = 428\,592$ and $433\,945$), each sitting on the constraint
($\mathrm{CVaR}_{5\%}\approx 0$) as the solution of a chance-constrained program
must, in $1.67$\,h and $4.14$\,h of one rented GPU. Batching an entire compass
sweep is worth $47\times$ over the NumPy reference; compiling the Krylov
iteration into a single call is worth a further $2.2$--$2.8\times$ at the larger
geometry, with the objective unchanged to $10^{-7}$.

\paragraph{The certification question.} A policy is reported only after being
re-solved on scenarios no run has seen, on the independent NumPy reference. The
accepted policy gives $391/400$ compliant scenarios on a fresh draw,
reliability $0.978$ with a Wilson interval of $[0.958, 0.988]$, clearing the
$0.94$ gate; its cost reproduces to $0.02\%$. The same protocol rejects a policy
optimized on sixteen scenarios: it reported reliability $1.000$ on its own
sample and achieves $0.912$ $[0.881, 0.936]$ on four hundred unseen ones.

\subsection{The methodological result}

The finding that generalizes beyond this problem is a statement about which
statistics of a stochastic program transport from the optimization sample to the
aquifer.

For the accepted policy on three disjoint draws of $400$ scenarios, the cost
reproduces to $0.02\%$ and the reliability to $0.003$
($0.978$, $0.975$, $0.978$), while the CVaR tail takes the values $+0.616$,
$+0.254$ and $+0.334$. Those three are not distinguishable: their spread is
smaller than the standard error of any one of them ($\pm0.21$ to $\pm0.47$),
because a $5\%$ tail of $400$ draws averages twenty of the most variable
realizations in the set. The reported in-sample value, $+0.005$, lies outside all
three intervals --- it is a fourth sample, not a better estimate.

The practical consequences, which are cheap to adopt and are the recommendation
of this thesis, are fourfold. Certify a policy with \emph{reliability and its
interval}, not with a tail value. Quote a tail \emph{with its standard error}, and
only as evidence about where the policy sits relative to the constraint. Do not
rank policies by tail value at sample sizes that are affordable: at $M=400$ the
tail cannot separate policies that reliability separates easily. And re-evaluate
on draws that no part of the optimization has touched, because scenarios that
continue the training seed reproduce the run rather than test it.

\subsection{Limitations of the present evidence}

\begin{itemize}
  \item \textbf{One two-dimensional geometry.} All 2D conclusions concern a single
        $200\times120$\,m configuration, nine wells and one source. The
        observation that well rates are seed-insensitive --- two independently
        seeded searches followed a digit-for-digit identical trajectory --- is a
        property of this cost model, not a general law, and it means the reported
        policy is one of several near-equivalent designs rather than a unique
        optimum.
  \item \textbf{Mixed precision at the margin.} The accepted policy was found in
        mixed float32 (float32 operator, float64 Krylov recurrence). That mode is
        not unconditionally safe: two of three fully float32 runs aborted on a
        non-finite transport residual, which is why the search now tolerates a
        single rejected candidate but not an entire diverging sweep. The
        float32/float64 policy gap of $1.2\%$ is within the noise that fine
        compass steps introduce.
  \item \textbf{Tail resolution.} At $M=400$ the tail statistic is not resolvable
        to better than a few tenths. Establishing the tail of the distribution to
        the precision at which reliability is known would require thousands of
        scenarios, i.e.\ an order of magnitude more transport solves than the
        optimization itself.
  \item \textbf{The three-dimensional model is verified but not certified, and
        is coarse.} Seven gates are asserted on the solver --- exact seepage
        velocity, transversally invariant head, a Thiem shape to $8.1\times
        10^{-3}$, agreement with the 2D code under extrusion to $10^{-13}$, the
        advection front speed to $1.2\%$, a breakthrough variance consistent
        with $\alpha_L v$, and a mass ledger closing to $10^{-15}$ --- but there is
        no analytical \emph{shape} gate for 3D transport, deliberately: the
        Ogata--Banks solution is a continuous Dirichlet step, and against a finite
        inlet it differs by a truncation error that shrinks with time ($0.32\to
        0.13$), so it would test the upwinding both codes share rather than the
        dimensional extension. What the 3D results lack is a \emph{certificate}:
        the paper preset's interval is $[0.922,0.988]$ and the $8\,000$-cell
        policy's is $[0.872,0.963]$, neither clearing the $0.94$ gate, and the
        4\,096-cell geometry is an order of magnitude coarser than the 2D
        campaign with a single log-normal conductivity family rather than a
        layered geology.
  \item \textbf{Timings are single-machine.} Reported GPU costs come from rented
        T4 and A100 instances with the batch sizes stated; the fixed-cost model
        behind the larger projections is a two-point fit that was not
        reproducible to better than about $15\%$.
\end{itemize}

\subsection{Work that would remove these limitations}

\begin{enumerate}
  \item \textbf{A certified 3D policy, and an affordable way to compute one.}
        Certifying at the 2D standard means $M_{\text{test}}\approx400$ and
        $M_{\text{train}}\approx100$ --- about $18$\,h per two-start search on
        one CPU at $8\,000$ cells --- so the work is as much computational as
        statistical: the batched torch backend on the $32\,768$-cell preset, and
        a search that uses the adjoint instead of a thousand forward solves.
  \item \textbf{A 2.5D axisymmetric bridge} between the 2D and 3D models, which
        would give a second independent route to the same answer and make the
        dimensional reduction a check instead of an assumption.
  \item \textbf{Operator learning in 3D} as an input--output map between
        conductivity realizations and compliance-point histories, with the same
        seed-sweep and resolution-transfer protocol applied in 1D, so that the
        surrogate claim is either established or refuted at the resolution the
        optimization actually uses.
  \item \textbf{Tail estimation at scale.} If a tail statement is wanted, it has
        to be paid for: $M$ of a few thousand, drawn from a frozen policy, with
        the standard error reported as a matter of course. Until then the honest
        claim is about reliability.
  \item \textbf{Adaptive design, solved by progressive hedging.} The formulation
        is written down and checked --- first-stage rates, recourse rates that may
        depend on what has been observed at $t_{\text{obs}}$, nonanticipativity on
        the first stage only, and a value of the stochastic solution and of perfect
        information that a here-and-now model cannot report --- and the
        decomposition on the one-stage problem is explicitly degenerate (a
        nine-dimensional consensus block has nothing to decompose). Running it at
        benchmark scale is the experiment that would say whether observing an
        aquifer for a year is cheaper than installing more wells, which is the
        question the 3D cost wall makes urgent.
  \item \textbf{Removing the per-step kernel launches.} The measured $2.2$--$2.8\times$
        came from compiling the iteration body; the sequential dependence of
        implicit time stepping remains, and a formulation that solves the whole
        time history at once would remove the remaining synchronizations. Any
        such change must still pass the mass-balance and field-parity gates.
\end{enumerate}

The wider lesson is that the two questions the thesis opened with --- is the
forward model right, and is the optimizer's evidence about the aquifer or about
its own sample --- are answered by the same kind of device: an independent
computation that can disagree. The cost of each device is small relative to the
optimization it audits, and in this work every one of them changed a conclusion.
